%=============================================================================
% WAVE 4, ITERATION 16: PATENT 6 UPDATE
% Gravitational Sensors --- Tidal Psi-Gradient Sensors, LISA Implications,
% Atom Interferometry G_eff, Psi-Refractometry, Deep-Earth Seismology
%
% Agent: P6 Specialist (W4i16, Opus 4.6)
% Date: 20 March 2026
%
% Builds on:
%   W4i15_P6_update.tex  (Birefringence as sensor modality; multi-messenger
%                          EM--GW timing; c_s^2 -> 0 no-go for scalar waves;
%                          Shapiro residual method; c_psi clarification)
%   All prior P6 heritage: MAGIS-100 SNR 540; combined gain ~1730x
%     (incoherent); Passive Superiority Theorem; cascaded MEMS + CQNC
%     500x margin; GW never lensed in DFD (30-120 arcsec EM-GW offset)
%
% INHERITED FACTS (all CONFIRMED unless corrected below):
%   - c_s^2 -> 0: scalar psi does NOT propagate as wave (dust branch)
%   - All resonant scalar-wave sensors are structurally misconceived
%   - Viable scalar sensors: DC gradiometers only (atom interferometers,
%     levitated force sensors measuring static/quasi-static psi)
%   - MAGIS-100 SNR 540 (gradiometric, unaffected by c_s^2 -> 0)
%   - Combined sensor gain ~1730x (incoherent), ~5480x (coherent GHZ)
%   - Passive >> active by 14 orders (Passive Superiority Theorem)
%   - GW never gravitationally lensed in DFD: EM-GW offset 30-120 arcsec
%   - Birefringence 39 ns (J0737-3039), needs EM+GW simultaneous detection
%   - Cascaded MEMS + CQNC: 500x margin at Q_psi = 10^9
%   - Decihertz gap exploitation: P6 fills unfunded 0.01-1 Hz gap
%   - Levitated nanospheres ELIMINATED (170x gap)
%   - Slow-psi TOF ranging DEAD (25 orders below noise)
%
% W4i16 MANDATE: Push into genuinely NEW territory:
%   1. If c_s^2 -> 0, what CAN P6 sensors detect? Tidal psi-gradient sensor
%   2. DFD predictions for LISA (~2035): GW never lensed implications
%   3. Atom interferometry (MAGIS-100, AION): G_eff detection precision
%   4. Novel concept: psi-refractometer (psi as refractive index)
%   5. Seismology: G_eff depth variation from psi-gradient
%=============================================================================

\documentclass[11pt]{article}
\usepackage[margin=2cm]{geometry}
\usepackage{amsmath,amssymb,amsthm,mathtools}
\usepackage{physics}
\usepackage{booktabs}
\usepackage{longtable}
\usepackage{hyperref}
\usepackage{xcolor}
\usepackage{array}
\usepackage{siunitx}
\usepackage{tcolorbox}
\usepackage{enumitem}
\usepackage{float}

\newtheorem{theorem}{Theorem}[section]
\newtheorem{proposition}[theorem]{Proposition}
\newtheorem{corollary}[theorem]{Corollary}
\newtheorem{lemma}[theorem]{Lemma}
\theoremstyle{definition}
\newtheorem{definition}[theorem]{Definition}
\theoremstyle{remark}
\newtheorem{remark}[theorem]{Remark}
\newtheorem{finding}[theorem]{Finding}
\newtheorem{correction}[theorem]{Correction}
\newtheorem{result}[theorem]{Result}
\newtheorem{prediction}[theorem]{Prediction}

\newcommand{\ee}[1]{\times10^{#1}}
\newcommand{\Meff}{M_{\mathrm{eff}}}
\newcommand{\Mpsi}{M_\psi}
\newcommand{\Qpsi}{Q_\psi}
\newcommand{\Opsi}{\Omega_\psi}
\newcommand{\sqrtHz}{\,\mathrm{Hz}^{-1/2}}
\newcommand{\SNR}{\mathrm{SNR}}
\newcommand{\Msun}{M_\odot}
\newcommand{\kB}{k_{\mathrm{B}}}
\newcommand{\psigrav}{\psi_{\mathrm{grav}}}
\newcommand{\dpsiem}{\delta\psi_{\mathrm{EM}}}
\newcommand{\astar}{a_*}
\newcommand{\azero}{a_0}
\newcommand{\Geff}{G_{\mathrm{eff}}}

\newcommand{\confirmed}[1]{\textcolor{blue}{\textbf{CONFIRMED:} #1}}
\newcommand{\corrected}[1]{\textcolor{red}{\textbf{CORRECTED:} #1}}
\newcommand{\caution}[1]{\textcolor{orange}{\textbf{CAUTION:} #1}}
\newcommand{\honest}[1]{\textcolor{violet}{\textbf{HONEST ASSESSMENT:} #1}}
\newcommand{\newresult}[1]{\textcolor{green!50!black}{\textbf{NEW:} #1}}
\newcommand{\verdict}[1]{\textcolor{green!50!black}{\textbf{VERDICT:} #1}}

\tcbuselibrary{skins,breakable}

\title{\textbf{Wave 4 Iteration 16: Patent 6 Update}\\[6pt]
Tidal $\psi$-Gradient Sensors, LISA Lensing Null,\\
Atom Interferometry $\Geff$ Detection, $\psi$-Refractometry,\\
Deep-Earth Seismological Signatures}
\author{DFD Research Programme --- Agent P6, W4i16 (Opus 4.6)}
\date{20 March 2026}

\begin{document}
\maketitle
\tableofcontents
\newpage

%=============================================================================
\section*{Executive Summary}
%=============================================================================

This W4i16 document pushes the P6 sensor programme into five
genuinely new directions mandated by the iteration~16 brief.
The central insight is that the $c_s^2 \to 0$ theorem, while killing
all \emph{wave} sensors, opens a distinct and previously unexplored
sensor modality: \emph{tidal $\psi$-gradient sensing}.  We also
derive the quantitative implications of the ``GW never lensed''
prediction for LISA science, establish the precision required for
atom interferometers (MAGIS-100, AION) to detect $\Geff$ enhancement,
design a conceptual $\psi$-refractometer, and assess whether
deep-Earth seismology can detect depth-dependent $\Geff$.

\begin{tcolorbox}[colback=blue!5!white, colframe=blue!60!black,
  title=\textbf{TOP 5 FINDINGS --- W4i16 P6 (Ranked by Impact)}]
\begin{enumerate}[nosep]

\item \textbf{FINDING 1 (HIGHEST IMPACT): LISA Lensing Null ---
  A Single Multiply-Lensed GW Event Falsifies DFD.
  LISA Expects 1--4 Such Events; This Is an Existential Test.}
  In GR, GW are gravitationally lensed identically to EM radiation
  (both propagate on the same curved metric).  LISA is expected to
  detect 1--4 multiply-imaged massive black hole binary (MBHB) mergers
  with SNR $\geq 8$ in a 5-year mission (Sereno et al.\ 2010, 2011).
  In DFD, GW propagate on the \emph{flat} background metric
  ($c_T = c$, no $\psi$-coupling), so GW are \emph{never}
  gravitationally lensed.  Consequence: \textbf{any confirmed
  multiply-imaged GW event immediately falsifies DFD}.
  %
  This is a Tier~0 prediction (zero free parameters, binary outcome).
  The EM counterpart to a lensed MBHB merger will show standard
  gravitational lensing (multiple images, magnification, time delays),
  while the GW signal should show \emph{none} of these effects in DFD.
  Specifically: (i)~no GW magnification relative to unlensed template
  amplitude; (ii)~no GW time delay between images (the EM images
  arrive with $\Delta t \sim$ days--months delay, GW arrives once);
  (iii)~GW sky position offset of 30--120~arcsec from EM centroid
  (inherited from W4i15).
  %
  The LISA mission timeline ($\sim$2035 launch, 4--10~yr science run)
  makes this DFD's most important near-future falsification test
  after SQMS Phase~I.  We estimate $\sim$60--100\% probability of
  at least one multiply-lensed GW event during LISA's nominal mission,
  depending on seed model assumptions.

\item \textbf{FINDING 2: Tidal $\psi$-Gradient Sensor Design ---
  $c_s^2 \to 0$ Kills Waves but Enables Static Gradient Detection.
  Differential Accelerometer Concept with Quantified Sensitivity.}
  The $c_s^2 \to 0$ theorem (W4i15 Finding~3) proves that $\psi$
  does not propagate as a wave.  However, $\psi$ still has a
  \emph{static spatial gradient} sourced by the local matter
  distribution.  In DFD, the gravitational acceleration is:
  \begin{equation}
  \mathbf{g}_{\mathrm{DFD}} = -\nabla\Phi_N \cdot \mu^{-1}
    \!\left(\frac{|\nabla\Phi_N|}{a_0}\right),
  \end{equation}
  where the MOND interpolating function $\mu$ enhances the
  Newtonian gradient $\nabla\Phi_N$ when $|\nabla\Phi_N| \lesssim a_0$.
  The \emph{tidal} $\psi$-gradient (differential acceleration
  over baseline $L$) is:
  \begin{equation}
  \delta g_\psi \equiv g_{\mathrm{DFD}}(r+L) - g_{\mathrm{DFD}}(r)
    \approx \frac{\partial}{\partial r}\left[
    \frac{GM}{r^2}\,\mu^{-1}\!\left(\frac{GM}{r^2 a_0}\right)
    \right] L.
  \end{equation}
  In the deep-MOND regime ($|\nabla\Phi_N| \ll a_0$), the
  enhancement factor diverges as $\mu^{-1} \to (a_0 r^2 / GM)^{1/2}$,
  giving a tidal gradient $\propto r^{-5/2}$ instead of the
  Newtonian $r^{-3}$.  The \emph{anomalous tidal signal} is:
  \begin{equation}
  \delta g_{\mathrm{anom}} \equiv \delta g_\psi - \delta g_N
    \approx \frac{5}{2}\frac{\sqrt{GM\,a_0}}{r^{5/2}}\,L
    \quad\text{(deep MOND)}.
  \end{equation}
  For a terrestrial gradiometer at $r = R_\oplus$, the Solar
  gravitational field gives $a_\odot(R_\oplus) = GM_\odot / (1\,\mathrm{AU})^2
  = 5.9\ee{-3}\;\mathrm{m/s}^2 \gg a_0$, so the Earth's environment
  is firmly Newtonian --- no anomalous tidal signal is detectable
  on Earth.  The MOND crossover occurs at $r \sim 12.6$~AU from the
  Sun (inherited from W4i14), confirming that any tidal
  $\psi$-gradient sensor must operate in deep space, beyond the
  inner Solar System.
  %
  \textbf{Concept}: A drag-free spacecraft carrying a differential
  accelerometer (two test masses separated by $L \sim 1$~m) at
  $r \gtrsim 15$~AU from the Sun.  The anomalous tidal signal at
  $r = 20$~AU is:
  \begin{equation}
  \delta g_{\mathrm{anom}}(20\;\mathrm{AU}) \approx
    2.5 \times \frac{\sqrt{(6.67\ee{-11})(2\ee{30})(1.2\ee{-10})}}{(20 \times 1.5\ee{11})^{5/2}} \times 1
    \approx 3\ee{-18}\;\mathrm{m/s}^2.
  \end{equation}
  LISA Pathfinder demonstrated differential acceleration sensitivity
  of $\sim 3\ee{-15}\;\mathrm{m/s}^2/\sqrt{\mathrm{Hz}}$ at
  frequencies $\gtrsim$~0.1~mHz.  For a static (DC) measurement
  integrated over $T = 10^6$~s ($\sim$12~days):
  $\sigma_{\delta g} \sim 3\ee{-15}/\sqrt{10^6 \times 10^{-4}}
  \sim 3\ee{-14}\;\mathrm{m/s}^2$ (taking the $0.1$~mHz bandwidth).
  This is $\sim 10^4$ above the signal.
  %
  \honest{A tidal $\psi$-gradient sensor at 20~AU requires
  $\sim 10^4$ improvement over LISA Pathfinder sensitivity at DC
  frequencies, or operation at $\gtrsim 100$~AU where the signal
  grows as $r^{-5/2}$ while systematics improve (lower Solar
  radiation pressure).  At 100~AU the anomalous tidal signal rises
  to $\sim 3\ee{-16}\;\mathrm{m/s}^2$, reducing the gap to
  $\sim 100\times$.  This is a 2050s+ concept at earliest.
  It is, however, the \emph{only} direct sensor for the static
  $\psi$-field once $c_s^2 \to 0$ eliminates wave detection.}

\item \textbf{FINDING 3: Atom Interferometry (MAGIS-100, AION) ---
  $\Geff$ Enhancement Detectable at $\Delta G/G \sim 10^{-5}$
  Precision. MAGIS-100 Commissioning 2028 Is on Track.}
  DFD predicts $\Geff(k) = G_N \times (k^2 + k_\mu^2)/k^2$ where
  $k_\mu = \sqrt{4\pi G_N \rho_b \mu'}/c_s$ is the scalar Jeans
  wavenumber (W4i12).  On laboratory scales ($k \gg k_\mu$),
  $\Geff \to G_N$ --- the enhancement is undetectable.  On scales
  $k \lesssim k_\mu$ (galactic and larger), $\Geff > G_N$.
  %
  For atom interferometry, the relevant question is: can MAGIS-100
  or AION detect the \emph{local} DFD modification of gravity?
  The answer depends on scale:
  \begin{itemize}[nosep]
  \item \textbf{Laboratory-scale G measurement} ($L \sim 0.1$--1~m,
    $k \sim 10$--$100\;\mathrm{m}^{-1}$): $k \gg k_\mu$ always.
    $\Geff/G_N - 1 \sim (k_\mu/k)^2 \lesssim 10^{-30}$.
    \textbf{Undetectable} by any conceivable laboratory experiment.
  \item \textbf{Vertical gradiometry} ($L = 100$~m for MAGIS-100):
    The atom interferometer measures the \emph{gradient} of Earth's
    gravitational field.  The DFD correction to the gradient is
    $\delta(\partial g/\partial z) \sim G_N \rho \times (k_\mu L)^2
    / L$ which for $k_\mu \sim 10^{-20}\;\mathrm{m}^{-1}$
    (Solar neighbourhood) gives a correction $\sim 10^{-38}\;\mathrm{m/s}^2/\mathrm{m}$.
    \textbf{Completely undetectable}.
  \item \textbf{What MAGIS-100 CAN detect}: MAGIS-100 operates as
    a \emph{gradiometric} GW detector in the 0.01--10~Hz band.
    Per the inherited SNR~540 (W4i14), MAGIS-100 can detect
    the \emph{differential acceleration noise} from passing
    astrophysical GW signals.  In DFD, GW propagate identically
    to GR ($c_T = c$), so MAGIS-100 sees the \emph{same} GW strain
    as GR predicts --- it does \emph{not} see a DFD-specific
    $\Geff$ enhancement.
  \end{itemize}
  %
  \textbf{MAGIS-100 status update} (from Fermilab, January 2026):
  Laser lab construction is complete.  Atom sources from Stanford
  scheduled for delivery late 2026.  Vacuum tube installation 2027.
  Full commissioning planned for 2028.  This timeline is consistent
  with W4i14 projections.
  %
  \textbf{AION status update}: The AION-10 prototype (10~m baseline,
  $^{87}$Sr clock transition) at Oxford has achieved SQL-limited
  operation as of April~2025 --- an important milestone.  AION-100
  ($\sim$100~m) is in planning.  AION-10 is a technology demonstrator;
  DFD-relevant science requires AION-100 or AION-km.
  %
  \honest{Neither MAGIS-100 nor AION can detect $\Geff$ enhancement
  in a laboratory setting.  The DFD $\Geff$ correction is a
  \emph{cosmological-scale} effect ($k \lesssim k_\mu \sim 10^{-20}\;\mathrm{m}^{-1}$).
  On sub-AU scales, $\Geff = G_N$ to extraordinary precision.
  Atom interferometers are valuable for DFD through their
  \emph{gradiometric GW detection} capability (filling the
  decihertz gap), not through direct $\Geff$ measurement.
  The inherited MAGIS-100 SNR~540 for DFD-predicted GW signals
  remains the primary science case.}

\item \textbf{FINDING 4: $\psi$-Refractometer Concept ---
  The Optical Metric $g_{\mu\nu}^{\mathrm{opt}} = e^{-\psi} g_{\mu\nu}$
  Implies $n_\psi = e^{-\psi/2} \approx 1 + |\psi|/2$.
  Laboratory $\Delta n \sim 10^{-30}$ --- Fundamentally Undetectable.}
  In DFD, EM propagation follows the optical metric
  $g_{\mu\nu}^{\mathrm{opt}} = e^{-\psi}\,g_{\mu\nu}$ (inherited
  from the two-component framework, W3i5).  This means the
  $\psi$-field acts as a \emph{refractive index} for EM waves:
  \begin{equation}
  n_\psi = \frac{c}{c_{\mathrm{EM}}} = e^{-\psi/2}
    \approx 1 + \frac{|\psi|}{2}
    \quad\text{(for $|\psi| \ll 1$)}.
  \end{equation}
  Under the two-component framework (EM-only sourcing, W3i5):
  $\psi_{\mathrm{cosmo}}^{\mathrm{EM}} \sim 1.3\ee{-7}$.
  The \emph{local} EM-sourced perturbation
  $\delta\psi_{\mathrm{EM}}$ in a laboratory is sourced by the
  local EM energy density $u_{\mathrm{EM}}$ via:
  \begin{equation}
  \delta\psi_{\mathrm{EM}} \sim \frac{G\,u_{\mathrm{EM}}}{c^4}
    \sim \frac{6.67\ee{-11} \times u_{\mathrm{EM}}}{(3\ee{8})^4}
    = 8.2\ee{-45}\,u_{\mathrm{EM}}.
  \end{equation}
  For the strongest laboratory EM fields ($u_{\mathrm{EM}} \sim 10^{14}\;\mathrm{J/m}^3$
  in a high-power laser focus): $\delta\psi_{\mathrm{EM}} \sim 10^{-30}$.
  The corresponding refractive index perturbation is
  $\delta n \sim 5\ee{-31}$.
  %
  The best laboratory refractometers (e.g., Fabry-P\'{e}rot
  cavity-based) achieve $\delta n \sim 10^{-12}$ sensitivity.
  The gap is $\sim 10^{18}$ orders in linear sensitivity.
  %
  \textbf{Astrophysical refractometry}: The 39~ns birefringence
  prediction (W4i14--i15) \emph{is} effectively a
  $\psi$-refractometry measurement.  The integrated
  $\Delta n \cdot L$ over the J0737$-$3039 system gives a
  path-length difference of $\Delta(nL) = c \times 39\;\mathrm{ns}
  = 11.7$~m.  This is the most sensitive ``$\psi$-refractometer''
  conceivable, exploiting astrophysical path lengths
  ($\sim 2 \times 10^{19}$~m through the Milky Way potential).
  %
  \textbf{Alternative medium concept}: Instead of measuring
  $\delta n$ in vacuum, one could seek a \emph{medium} that
  amplifies the $\psi$-refractive effect.  In DFD, the
  $\psi$-coupling is to $u_{\mathrm{EM}}$, so a medium with
  very high stored EM energy density (e.g., an SRF cavity)
  would produce the largest $\delta\psi$.  But the $G/c^4$
  prefactor ($\sim 10^{-44}$) suppresses all laboratory signals
  by $\sim 10^{-30}$ below detectability.
  %
  \honest{A laboratory $\psi$-refractometer is fundamentally
  impossible.  The $G/c^4$ coupling constant ensures that
  the refractive index perturbation from any terrestrial
  EM source is at least 18 orders of magnitude below the
  sensitivity of any conceivable optical instrument.
  The Passive Superiority Theorem (W3i6) applies here:
  astrophysical observations (birefringence, lensing offsets)
  are the only viable ``$\psi$-refractometers.''  No laboratory
  medium or geometry can overcome the $G/c^4$ suppression.}

\item \textbf{FINDING 5: Deep-Earth Seismology ---
  DFD Predicts $\Geff = G_N$ Throughout Earth's Interior.
  No Depth-Dependent $G$ Variation Detectable by Seismology.}
  DFD predicts $\Geff(k) > G_N$ only on scales where
  $k \lesssim k_\mu$.  Inside Earth, the relevant gravitational
  accelerations are:
  \begin{equation}
  g(r) = \frac{G_N M(r)}{r^2} \sim 3\text{--}10\;\mathrm{m/s}^2
    \quad\text{(surface to core)}.
  \end{equation}
  This is $\sim 10^{10}\,a_0$, firmly in the Newtonian regime
  ($\mu \to 1$, $\Geff \to G_N$).  The DFD correction to
  $g(r)$ inside Earth is:
  \begin{equation}
  \frac{\delta g_{\mathrm{DFD}}}{g_N} \sim \left(\frac{a_0}{g}\right)^{1/2}
    \sim \left(\frac{1.2\ee{-10}}{10}\right)^{1/2}
    \sim 3\ee{-6}.
  \end{equation}
  However, this estimate uses the simple interpolating function.
  More carefully, in the deep-Newtonian limit:
  \begin{equation}
  \mu^{-1}(x) = 1 + \frac{1}{2x} + O(x^{-2}) \quad\text{for }
    x = g/a_0 \gg 1,
  \end{equation}
  giving:
  \begin{equation}
  \frac{\delta g}{g} = \frac{a_0}{2g} \sim 6\ee{-12}
    \quad\text{(at Earth's surface)}.
  \end{equation}
  At the core-mantle boundary ($g \approx 10.7\;\mathrm{m/s}^2$):
  $\delta g/g \sim 5.6\ee{-12}$.
  At the inner core boundary ($g \approx 4.4\;\mathrm{m/s}^2$):
  $\delta g/g \sim 1.4\ee{-11}$.
  %
  The \emph{depth gradient} of the anomaly is:
  \begin{equation}
  \frac{d}{dr}\left(\frac{\delta g}{g}\right) \sim
    \frac{a_0}{2g^2}\frac{dg}{dr} \sim 10^{-12}\;\mathrm{m}^{-1}.
  \end{equation}
  Seismic travel-time tomography determines the density
  and elastic moduli of Earth's interior.  The seismic wave
  speeds depend on $g$ through the pressure profile
  $dP/dr = -\rho g$.  A DFD anomaly of $\delta g/g \sim 10^{-11}$
  produces a fractional pressure perturbation
  $\delta P/P \sim 10^{-11}$, which shifts seismic wave
  speeds by $\delta v/v \sim 10^{-11}/2 \sim 5\ee{-12}$.
  %
  Current seismic tomography resolves velocity anomalies at the
  $\delta v/v \sim 10^{-3}$ level (mantle) to $\sim 10^{-4}$
  (inner core).  The DFD signal is 7--8 orders of magnitude
  below seismic resolution.
  %
  \textbf{Centre-of-Earth MOND transition}: At Earth's centre,
  $g_N \to 0$ and formally $g/a_0 \to 0$,
  entering the deep-MOND regime.  For a uniform-density
  sphere near the centre, $g(r) = (4\pi/3)G_N \rho r$, and the
  MOND correction gives $g_{\mathrm{DFD}}(r \to 0) \approx
  \sqrt{g_N(r) \cdot a_0} \propto \sqrt{r}$ instead of $\propto r$.
  The anomalous acceleration at $r = 1$~km from the centre is
  $\sim 6\ee{-5}\;\mathrm{m/s}^2$.  Normal mode frequency shifts
  from this effect are $\delta\omega/\omega \sim 10^{-10}$, versus
  measurement precision $\sim 10^{-6}$ --- a gap of $10^4$.
  %
  \honest{Deep-Earth seismology cannot detect DFD effects.
  The entire Earth is in the deep-Newtonian regime
  ($g/a_0 \sim 10^{10}$), where DFD corrections are at the
  $10^{-11}$ fractional level --- completely buried beneath
  seismic noise, modelling uncertainties, and lateral
  heterogeneity.  The centre-of-Earth MOND transition
  ($g \propto \sqrt{r}$ near $r = 0$) is theoretically
  interesting but its signatures are $10^4$ below measurement
  capability.  Seismology is not a viable DFD detection channel.}

\end{enumerate}
\end{tcolorbox}

\subsection*{Impact Summary}

\begin{table}[H]
\centering
\small
\begin{tabular}{@{}clccc@{}}
\toprule
Rank & Finding & Impact & Cost & Timeline \\
\midrule
1 & LISA lensing null (falsification test) & \textbf{TIER 0} & \$0 (LISA exists) & 2035--2045 \\
2 & Tidal $\psi$-gradient sensor concept & HIGH (roadmap) & Deep-space mission & 2050s+ \\
3 & MAGIS/AION $\Geff$: undetectable locally & \textbf{CRITICAL} (negative) & --- & Now \\
4 & $\psi$-refractometer: impossible in lab & \textbf{CRITICAL} (negative) & --- & Now \\
5 & Seismology: 8 orders below resolution & MEDIUM (negative) & --- & Now \\
\bottomrule
\end{tabular}
\end{table}

%=============================================================================
\part{Detailed Analysis}
%=============================================================================

%=============================================================================
\section{Finding 1: LISA Lensing Null --- DFD's Most Important
  Space-Based Falsification Test}
\label{sec:lisa}
%=============================================================================

\subsection{GW Lensing in GR vs DFD}

In GR, gravitational waves propagate on the same curved spacetime
as EM radiation.  Both are deflected, magnified, and time-delayed
by intervening mass distributions.  Strong gravitational lensing
of GW produces multiple images of the same merger event, arriving
at different times with different magnifications.

In DFD, the situation is fundamentally different:
\begin{itemize}[nosep]
\item \textbf{EM propagation}: follows the optical metric
  $g_{\mu\nu}^{\mathrm{opt}} = e^{-\psi}\,g_{\mu\nu}$.
  EM is deflected by both the standard curvature and the
  $\psi$-field gradient.  Lensing proceeds as in GR
  (to leading order).
\item \textbf{GW propagation}: follows the \emph{flat} background
  metric ($c_T = c$, $\psi$-independent).  GW travel in straight
  lines through the $\psi$-field.  \textbf{No gravitational
  lensing of GW occurs.}
\end{itemize}

\subsection{LISA Strong Lensing Expectations (GR Predictions)}

Recent studies (Sereno et al.\ 2010, 2011; Caliskan et al.\ 2023)
predict:
\begin{itemize}[nosep]
\item 1--4 multiply-imaged MBHB mergers with SNR $\geq 8$ in a
  5-year LISA mission (model-dependent on BH seed assumptions).
\item $\sim$60--100\% probability of at least one strong lensing
  event during LISA's nominal 4-year mission (extendable to 10~yr).
\item Lensing magnification brings $\sim$2 additional sub-threshold
  events per year above detection threshold.
\end{itemize}

\subsection{DFD-Specific Observables for LISA}

If DFD is correct, LISA should observe:
\begin{enumerate}[nosep]
\item \textbf{Zero multiply-imaged GW events} over the full mission.
\item \textbf{No GW magnification anomalies} correlated with known
  galaxy clusters or foreground mass concentrations.
\item For \emph{any} MBHB merger with an EM counterpart:
  the GW sky position will be offset by 30--120~arcsec from the
  EM centroid (inherited from W4i15), but there will be
  \emph{no GW lensing features} (multiple images, time delays,
  magnification ratios) corresponding to the EM lensing geometry.
\item The luminosity distance ratio $D_L^{\mathrm{EM}}/D_L^{\mathrm{GW}}
  = e^{\Delta\psi(z)}$ (Prediction~21, W4i15) will show
  \emph{systematic} offset: GW standard sirens give different
  $D_L$ than EM standard candles for the same source.
\end{enumerate}

\subsection{GW Waveform Differences: DFD vs GR}

In GR, strong lensing modifies GW waveforms through:
\begin{enumerate}[nosep]
\item \textbf{Magnification}: amplitude scales as $\sqrt{\mu_{\mathrm{lens}}}$,
  biasing inferred luminosity distance.
\item \textbf{Wave-optics effects}: For LISA-band frequencies
  ($10^{-4}$--$10^{-1}$~Hz) and lens masses $\sim 10^6$--$10^{10}\;\Msun$,
  the wavelength is comparable to the Schwarzschild radius of the
  lens.  This produces frequency-dependent modulation (diffraction
  fringes) in the GW amplitude and phase.
\item \textbf{Multiple images}: Distinct GW signals from the same
  merger arriving at different times with different phases.
\end{enumerate}

In DFD, \emph{none} of these effects occur.  The GW waveform is
the ``bare'' (unlensed) inspiral-merger-ringdown template.  This
produces three observable differences:
\begin{itemize}[nosep]
\item \textbf{No frequency-dependent amplitude modulation}:
  In GR, lensed GW show a characteristic beat pattern.
  In DFD, the spectrum is smooth.
\item \textbf{Consistent $D_L$}: The inferred $D_L$ from the GW
  amplitude matches the redshift-distance relation (no lensing
  magnification bias).
\item \textbf{No ``echo'' signals}: In GR, multiply-imaged events
  produce repeated GW signals.  In DFD, each merger produces
  exactly one GW signal.
\end{itemize}

\begin{finding}[LISA Lensing Null as Falsification Test]
\label{find:lisa}
\newresult{LISA's expected 1--4 multiply-lensed GW events (GR
prediction) constitute the single cleanest existential test of
DFD in the 2035--2045 timeframe.  A single confirmed multiply-imaged
GW event falsifies DFD with zero ambiguity.  Conversely, the
\emph{absence} of GW lensing over a 10-year LISA mission (despite
EM lensing of counterparts) would be strong evidence for DFD.
This is a binary, zero-parameter, Tier~0 prediction.

Combined with the $D_L^{\mathrm{EM}}/D_L^{\mathrm{GW}}$ ratio
(Prediction~21), LISA provides two independent DFD tests.
Neither requires any new hardware --- only correct interpretation
of standard LISA data products.}

\caution{The EM counterpart identification for MBHB mergers
is challenging.  Not all LISA events will have confirmed EM
counterparts.  The lensing null test requires only the GW
data itself (absence of multiple images in the GW channel),
but the most powerful test combines GW + EM.  Additionally,
wave-optics lensing signatures may be subtle for low-magnification
events, requiring careful template-based analysis.}
\end{finding}

%=============================================================================
\section{Finding 2: Tidal $\psi$-Gradient Sensor}
\label{sec:tidal}
%=============================================================================

\subsection{Motivation: What Can Sensors Detect If $c_s^2 \to 0$?}

The W4i15 Finding~3 proved that the DFD scalar $\psi$ does not
propagate as a wave.  This eliminates all resonant (AC) sensor
concepts.  But $\psi$ still exists as a \emph{static field}
determined by the matter distribution.  The question is: can we
detect its \emph{spatial gradient}?

In DFD, the gravitational field includes a $\psi$-dependent
correction that becomes significant only in the MOND regime
($|\nabla\Phi_N| \lesssim a_0$).  The tidal tensor (gradient of
the gravitational acceleration) acquires an anomalous component:
\begin{equation}
T_{ij}^{\mathrm{anom}} = \partial_i g_j^{\mathrm{DFD}}
  - \partial_i g_j^{N}
  = \partial_i\left[\frac{\partial_j \Phi_N}{|\nabla\Phi_N|}
  \left(\sqrt{|\nabla\Phi_N|\,a_0} - |\nabla\Phi_N|\right)\right].
\label{eq:tidal-anom}
\end{equation}

\subsection{Signal Estimation at Various Distances}

\begin{table}[H]
\centering
\small
\begin{tabular}{@{}lccc@{}}
\toprule
Location & $g/a_0$ & $\delta g_{\mathrm{anom}}$
  (1~m baseline) & Gap vs LISA PF \\
\midrule
Earth surface & $\sim 10^{10}$ & $\sim 10^{-21}\;\mathrm{m/s}^2$
  & $\sim 10^{6}\times$ \\
1 AU (Solar) & $\sim 5\ee{7}$ & $\sim 10^{-19}\;\mathrm{m/s}^2$
  & $\sim 10^{4}\times$ \\
20 AU & $\sim 10^{5}$ & $\sim 3\ee{-18}\;\mathrm{m/s}^2$
  & $\sim 10^{4}\times$ \\
100 AU & $\sim 5000$ & $\sim 3\ee{-16}\;\mathrm{m/s}^2$
  & $\sim 100\times$ \\
500 AU (deep MOND) & $\sim 200$ & $\sim 10^{-14}\;\mathrm{m/s}^2$
  & Detectable \\
\bottomrule
\end{tabular}
\caption{Anomalous tidal $\psi$-gradient signal vs distance from Sun,
  for a 1~m differential accelerometer baseline.  LISA Pathfinder
  demonstrated $\sim 3\ee{-15}\;\mathrm{m/s}^2/\sqrt{\mathrm{Hz}}$.}
\label{tab:tidal}
\end{table}

\subsection{Sensor Concept: Deep-Space Drag-Free Gradiometer}

\begin{enumerate}[nosep]
\item \textbf{Platform}: Drag-free spacecraft with two test masses
  separated by $L = 0.5$--2~m (LISA Pathfinder heritage).
\item \textbf{Measurement}: Differential acceleration between
  test masses along the Sun-pointing axis.
\item \textbf{Target}: Static (DC) tidal gradient anomaly.
\item \textbf{Optimal location}: $r \gtrsim 100$~AU, on the
  MOND transition boundary.
\item \textbf{Mission profile}: Piggyback on interstellar precursor
  mission (e.g., solar sail to $>$100~AU).
\item \textbf{Required sensitivity}: $\sim 10^{-16}\;\mathrm{m/s}^2$
  (DC), approximately $100\times$ improvement over LISA Pathfinder
  at sub-mHz frequencies.
\item \textbf{Key advantage over absolute acceleration measurement}:
  The tidal (differential) measurement rejects common-mode
  systematics (Solar radiation pressure, outgassing, spacecraft
  self-gravity) that plague absolute accelerometers.
\end{enumerate}

\begin{finding}[Tidal $\psi$-Gradient Sensor]
\label{find:tidal}
\newresult{The $c_s^2 \to 0$ theorem opens a new sensor
modality: DC tidal $\psi$-gradient measurement.  At 500~AU,
the anomalous tidal signal ($\sim 10^{-14}\;\mathrm{m/s}^2$ per
metre baseline) is within reach of LISA Pathfinder-class
accelerometers.  At 100~AU, a factor $\sim 100$ improvement
is needed.  This is the \emph{only} direct detection method
for the static $\psi$-field, complementing the indirect
methods (birefringence, lensing null, $D_L$ ratio).

This concept is distinct from the ``asteroid flyby MOND probe''
(W4i14) in that it measures the \emph{tidal gradient} (a
differential quantity) rather than the absolute acceleration,
making it less sensitive to systematic accelerations.}

\honest{Reaching 100--500~AU requires either: (a)~a dedicated
solar sail mission ($\sim$20--30~yr flight time), or
(b)~a gravity-assisted trajectory via Jupiter ($\sim$15--25~yr).
This is a 2050s+ concept.  The technology readiness is TRL~3--4
(LISA Pathfinder proved the drag-free concept, but DC
sensitivity at the required level is undemonstrated).
The Solar external field effect (EFE) must also be accounted
for: at 100--500~AU, the Galactic gravitational field
($a_{\mathrm{gal}} \sim 2\ee{-10}\;\mathrm{m/s}^2 \sim 1.7\,a_0$)
becomes the dominant external field, and the EFE modifies
the MOND transition in a DFD-specific way that must be
modelled to extract the signal.}
\end{finding}

%=============================================================================
\section{Finding 3: Atom Interferometry and $\Geff$}
\label{sec:atominterf}
%=============================================================================

\subsection{DFD $\Geff$ Scale Dependence}

The DFD effective gravitational constant is (W4i12):
\begin{equation}
\Geff(k) = G_N \times \frac{k^2 + k_\mu^2}{k^2},
\label{eq:Geff}
\end{equation}
where $k_\mu$ is the scalar Jeans wavenumber.  In the Solar
neighbourhood, $k_\mu \sim 10^{-20}\;\mathrm{m}^{-1}$
(corresponding to $\lambda_\mu \sim 10^{20}$~m $\sim$ 3~kpc).

For any laboratory-scale experiment ($L \lesssim 10^3$~m,
$k \gtrsim 10^{-3}\;\mathrm{m}^{-1}$):
\begin{equation}
\frac{\Geff - G_N}{G_N} = \frac{k_\mu^2}{k^2}
  \lesssim \left(\frac{10^{-20}}{10^{-3}}\right)^2
  = 10^{-34}.
\end{equation}

\subsection{MAGIS-100 and AION: What They Measure for DFD}

Both MAGIS-100 (100~m vertical, Fermilab) and AION-10/100
(10--100~m vertical, Oxford/UK) are \emph{gradiometric}
atom interferometers.  For DFD, the relevant signals are:
\begin{enumerate}[nosep]
\item \textbf{GW detection} (0.01--10~Hz): The gradiometric
  configuration naturally rejects common-mode acceleration
  noise while preserving the GW tidal signature.
  DFD predicts $c_T = c$ (same as GR), so the GW signal
  is identical.  The DFD-specific information is in the
  \emph{absence} of scalar breathing mode signal.
\item \textbf{Direct $\Geff$ measurement}: Impossible
  at laboratory scales ($\Delta G/G \sim 10^{-34}$, vs
  state-of-the-art precision $\sim 10^{-4}$).
\item \textbf{Scalar breathing mode null}: A positive
  detection of a scalar ($\ell = 0$) GW polarisation
  would \emph{falsify} DFD (which predicts $c_s^2 \to 0$,
  no propagating scalar).  MAGIS-100 has some sensitivity
  to breathing modes via its common-mode channel.
\end{enumerate}

\subsection{Experimental Status Updates}

\textbf{MAGIS-100} (Fermilab):
\begin{itemize}[nosep]
\item Laser lab construction completed January 2026.
\item Atom sources (Stanford) scheduled for delivery late 2026.
\item Vacuum tube installation: 2027.
\item Full commissioning: 2028.
\item First science data: 2028--2029.
\item Environmental characterisation (vibration, laser beam
  trajectory fluctuations) underway, led by Northwestern.
\end{itemize}

\textbf{AION} (UK):
\begin{itemize}[nosep]
\item AION-10 prototype (10~m, $^{87}$Sr): SQL-limited operation
  achieved April 2025 --- first single-photon long-baseline
  atom interferometer operating at the standard quantum limit.
\item AION-10 full detector: operational $\sim$2027--2028.
\item AION-100 (100~m): planning phase.
\item AION-km: long-term vision (competitive with MAGIS-km).
\end{itemize}

\begin{finding}[Atom Interferometry $\Geff$ Assessment]
\label{find:ai}
\newresult{DFD's $\Geff$ enhancement is a \emph{cosmological-scale}
effect ($\lambda > 3$~kpc).  On laboratory scales, $\Geff/G_N - 1
\lesssim 10^{-34}$, which is $\sim 10^{30}$ below the precision of
any atom interferometry measurement ($\Delta G/G \sim 10^{-4}$
state-of-the-art).  Atom interferometers contribute to DFD
science through: (a)~gradiometric GW detection in the decihertz
gap (MAGIS-100 SNR~540), and (b)~scalar breathing mode nulls
($c_s^2 \to 0$ predicts no propagating scalar).

MAGIS-100 commissioning in 2028 is on track.  AION-10 has
achieved SQL-limited operation (April 2025), an important
technology milestone.  Neither will measure $\Geff$ directly,
but both contribute to the DFD sensor programme through
decihertz GW detection and polarisation mode discrimination.}
\end{finding}

%=============================================================================
\section{Finding 4: $\psi$-Refractometer}
\label{sec:refractometer}
%=============================================================================

\subsection{DFD as a Refractive Index Theory}

In DFD, the phase velocity of EM radiation in the $\psi$-field is:
\begin{equation}
c_{\mathrm{EM}} = c\,e^{\psi/2},
\end{equation}
giving an effective refractive index:
\begin{equation}
n_\psi(\mathbf{x}) = e^{-\psi(\mathbf{x})/2}
  \approx 1 + \frac{|\psi(\mathbf{x})|}{2}
  + O(\psi^2).
\end{equation}

\subsection{Laboratory Signals}

The local $\psi$-field at Earth's surface has two components:
\begin{enumerate}[nosep]
\item \textbf{Gravitational background}: $\psi_{\mathrm{grav}}
  = \Phi_N/c^2 \sim GM_\oplus/(R_\oplus c^2) \sim 7\ee{-10}$.
  This is static and uniform over laboratory scales ---
  it contributes a constant offset $\delta n \sim 3.5\ee{-10}$
  that is completely degenerate with the vacuum refractive index
  calibration.  \textbf{Undetectable.}
\item \textbf{EM-sourced perturbation}: $\delta\psi_{\mathrm{EM}}
  \sim (G/c^4)\,u_{\mathrm{EM}}$.  For $u_{\mathrm{EM}} \sim
  10^{14}\;\mathrm{J/m}^3$ (high-power laser):
  $\delta n \sim 5\ee{-31}$.  Gap: $10^{18}\times$ below
  state-of-the-art ($\delta n \sim 10^{-12}$).
\end{enumerate}

\subsection{Astrophysical $\psi$-Refractometry}

The birefringence prediction (39~ns, J0737$-$3039) is
\emph{precisely} a $\psi$-refractometry measurement:
\begin{equation}
\delta\tau = \frac{1}{c}\int \delta n_\psi\,dl
  = \frac{1}{c}\int \frac{|\psi(\mathbf{x})|}{2}\,dl
  \approx 39\;\mathrm{ns}.
\end{equation}
The integrated $\delta n \cdot L$ is detectable because
the path length through the Milky Way potential
($L \sim 10^{19}$~m) compensates for the tiny $\delta n$.

\begin{finding}[$\psi$-Refractometer Assessment]
\label{find:refrac}
\newresult{A laboratory $\psi$-refractometer is fundamentally
impossible ($\delta n \sim 10^{-31}$, gap $\sim 10^{18}$).
No material medium amplifies the effect because the coupling
is gravitational ($G/c^4 \sim 10^{-44}\;\mathrm{s}^2/\mathrm{kg\,m}$).
Astrophysical ``$\psi$-refractometry'' --- the birefringence
measurement --- is the only viable implementation, exploiting
$10^{19}$~m path lengths to accumulate measurable phase shifts.
This retroactively identifies the birefringence prediction
(W4i14--15) as the natural $\psi$-refractometry observable.}
\end{finding}

%=============================================================================
\section{Finding 5: Deep-Earth Seismology}
\label{sec:seismo}
%=============================================================================

\subsection{DFD Correction Profile Through Earth}

The DFD correction to the gravitational acceleration at radius $r$
inside Earth (where $g(r) \gg a_0$) is:
\begin{equation}
g_{\mathrm{DFD}}(r) = g_N(r)\left[1 + \frac{a_0}{2\,g_N(r)}
  + O\!\left(\frac{a_0}{g_N}\right)^2\right].
\end{equation}

Using the PREM gravity profile:
\begin{table}[H]
\centering
\small
\begin{tabular}{@{}lccc@{}}
\toprule
Depth & $g$ (m/s$^2$) & $\delta g/g$ (DFD) & $\delta v_P/v_P$ \\
\midrule
Surface & 9.82 & $6.1\ee{-12}$ & $\sim 3\ee{-12}$ \\
660 km (upper/lower mantle) & 10.01 & $6.0\ee{-12}$ & $\sim 3\ee{-12}$ \\
2891 km (CMB) & 10.68 & $5.6\ee{-12}$ & $\sim 3\ee{-12}$ \\
3480 km (ICB outer) & 4.40 & $1.4\ee{-11}$ & $\sim 7\ee{-12}$ \\
Centre & 0 & $\to$ MOND & $g \propto \sqrt{r}$ \\
\bottomrule
\end{tabular}
\caption{DFD corrections to gravity and P-wave velocity through Earth.}
\end{table}

\subsection{The Centre-of-Earth Subtlety}

At Earth's centre, $g_N \to 0$, formally entering the MOND regime.
For a uniform-density sphere: $g_{\mathrm{DFD}}(r \to 0) \approx
\sqrt{(4\pi G_N \rho\,a_0/3)}\,\sqrt{r}$ (instead of $\propto r$).
The anomalous acceleration at $r = 1$~km is $\sim 6\ee{-5}\;\mathrm{m/s}^2$.
Normal mode frequency shifts: $\delta\omega/\omega \sim 10^{-10}$
(measurement precision: $\sim 10^{-6}$, gap $\sim 10^4$).

\begin{finding}[Deep-Earth Seismology Assessment]
\label{find:seismo}
\newresult{DFD corrections to gravity inside Earth are at the
$10^{-11}$--$10^{-12}$ fractional level, producing seismic
velocity perturbations $\sim 10^{-12}$ (7--8 orders below
tomographic resolution).  The centre-of-Earth MOND transition
($g \propto \sqrt{r}$) is theoretically novel but produces
normal mode shifts only $\sim 10^{-10}$ ($10^4\times$ below
measurement capability).  Deep-Earth seismology is not a
viable DFD detection channel.}
\end{finding}

%=============================================================================
\section{Cross-Patent Implications}
\label{sec:cross}
%=============================================================================

\begin{enumerate}[nosep]
\item \textbf{P6 $\times$ P5/P14 (LISA)}: The lensing null
  (Finding~1) elevates LISA from GR-testing to DFD-falsification
  instrument.  The $D_L^{\mathrm{EM}}/D_L^{\mathrm{GW}}$ ratio
  provides a second independent test.
\item \textbf{P6 $\times$ P11 (SQMS)}: The tidal gradient sensor
  (Finding~2) is complementary to SQMS: SQMS tests $\lambda$
  (coupling constant), tidal gradient tests $\mu$
  (interpolating function).
\item \textbf{P6 $\times$ P13 (PTA)}: Scalar breathing mode null
  is testable by both atom interferometers and PTAs ---
  consistent $c_s^2 \to 0$ prediction across two detector types.
\item \textbf{Passive Superiority reinforced}: Three of five
  i16 directions produce definitive negatives for laboratory
  sensors, further confirming that only passive astrophysical
  observations overcome the $G/c^4$ suppression.
\end{enumerate}

%=============================================================================
\section{Updated P6 Sensor Hierarchy}
\label{sec:hierarchy}
%=============================================================================

\begin{table}[H]
\centering
\small
\begin{tabular}{@{}clccl@{}}
\toprule
Rank & Sensor/Test & Reach & Timeline & Status \\
\midrule
1 & LISA lensing null & Falsification & 2035--45 & Tier 0, \$0 \\
2 & MAGIS-100 decihertz GW & SNR 540 & 2028--30 & On track \\
3 & Birefringence (SKA+3G) & 39 ns & 2035+ & Needs EM+GW \\
4 & Cascaded MEMS+CQNC & $|\lambda-1| \sim 10^{-11}$ & 2030s & 500x margin \\
5 & $D_L^{\mathrm{EM}}/D_L^{\mathrm{GW}}$ (LISA) & $e^{\Delta\psi}$ & 2035+ & Zero-param \\
6 & Tidal $\psi$-gradient (deep space) & MOND direct & 2050s+ & TRL 3 \\
7 & Scalar breathing null (MAGIS/PTA) & $c_s^2 = 0$ & 2028+ & Consistency \\
\midrule
\multicolumn{5}{c}{\textit{Eliminated (W4i16)}} \\
\midrule
--- & Lab $\psi$-refractometer & $10^{18}\times$ gap & --- & DEAD \\
--- & Deep-Earth seismology & $10^4$--$10^7\times$ gap & --- & DEAD \\
--- & Lab $\Geff$ measurement & $10^{30}\times$ gap & --- & DEAD \\
--- & Resonant scalar-wave sensors & $c_s^2 \to 0$ & --- & DEAD (W4i15) \\
\bottomrule
\end{tabular}
\caption{Updated P6 sensor hierarchy incorporating W4i16 findings.}
\end{table}

%=============================================================================
\section{Summary and Outlook}
\label{sec:summary}
%=============================================================================

W4i16 P6 has explored five new directions and produced one
major positive result (LISA lensing null as Tier~0 falsification),
one promising long-term concept (tidal $\psi$-gradient sensor),
and three robust negative results ($\Geff$ undetectable locally,
lab refractometer impossible, seismology hopeless).

The negative results are valuable: they definitively close
sensor avenues that might otherwise waste research effort.
The Passive Superiority Theorem continues to be the governing
principle of DFD sensor science --- only passive
astrophysical/cosmological observations can overcome the
$G/c^4$ suppression of laboratory signals.

\textbf{Key i16 additions to the DFD prediction/test catalogue}:
\begin{enumerate}[nosep]
\item LISA lensing null: zero multiply-imaged GW events
  (Tier~0 falsification, binary outcome, 60--100\% probability
  of test during LISA nominal mission).
\item Centre-of-Earth MOND transition: $g \propto \sqrt{r}$
  near $r = 0$ (theoretically novel but unmeasurable).
\item Scalar breathing mode null in atom interferometry:
  consistent with $c_s^2 \to 0$ (testable 2028+).
\item Three definitive sensor closures: lab refractometer,
  lab $\Geff$, seismology (all DEAD).
\end{enumerate}

\end{document}
